\section*{Transconduttanza}

\noindent\colorbox{orange}{
  \begin{minipage}{1.0\linewidth}
    \subsection*{Richiesta (a)}
    Nell'ipotesi che l'operazionale sia ideale e la sua uscita lineare, determinare
    un'espressione della transconduttanza dell'amplificatore $g_m=I_{OUT}/V_S$ in termini dei
    componenti circuitali. In particolare, verificare che:
    \begin{enumerate}
    \item[i.]
      come osservato al punto 2b, detta espressione non dipenda dal valore del carico;
    \item[ii.]
      sostituendo i valori numerici delle resistenze utilizzate, risulti un valore atteso per $g_m$ coincidente entro gli errori con quanto ottenuto al punto 1b.
    \end{enumerate}
  \end{minipage}
}
\subsection*{Risposta (a)}
Nell'ipotesi di idealit\'a degli op-amp e di operazione nel regime lineare \'e valido il \textit{principio di corto-circuito virtuale}
\begin{equation*}
  V_- = 0
\end{equation*}
applicando la KCL si ottiene la relazione fra $V_S$ e $I_{OUT}$ che permette di ottenere $g_m$
\begin{equation*}
  \begin{gathered}
    \frac{V_S}{R} = -\frac{V}{R} \\
    -\frac{V}{R} = -I_{OUT} + \frac{V}{R/10} \\ \implies
    \colorbox{blu}{\boxed{g_m = -\frac{11}{R}}}
  \end{gathered}
\end{equation*}

\noindent\colorbox{orange}{
  \begin{minipage}{1.0\linewidth}
    \subsection*{Richiesta (b)}
    Nelle stesse ipotesi di cui al punto precedente, determinare (a partire dalle rispettive
    definizioni) le impedenze di ingresso/uscita del circuito e verificare che siano compatibili
    con quanto misurato/dedotto ai punti 2a/2b.
  \end{minipage}
}
\subsection*{Risposta (b)}
Nell'ipotesi di idealit\'a dell'op-amp ed operazione in regime lineare, da quanto visto in precedenza
la corrente entrante per definizione
\begin{equation*}
  I_{in} = \frac{V_S - 0}{R}
\end{equation*}
di conseguenza l'impedenza d'ingresso
\begin{equation*}
  \colorbox{blu}{\boxed{Z_{in} \equiv \frac{V_S}{I_{in}} = R}}
\end{equation*}
Passivando il generatore di $V_S$ e aprendo il ramo di carico $R_L$, la corrente entrante da $V_{U1}$
deve essere uguale a quella uscente da $V$, tuttavia siccome
\begin{equation*}
  \frac{V - 0}{R} = \frac{0 - 0}{R} \implies V = 0
\end{equation*}
per la KCL vale
\begin{equation*}
  I_{U_1} = \frac{V - 0}{R} + \frac{V - 0}{R/10} \implies I_{U_1} = 0
\end{equation*}
di conseguenza
\begin{equation*}
  \colorbox{blu}{\boxed{Z_{out} \equiv \frac{V_t}{I_t} = \frac{V_{U_1} - V}{I_{U_1}} \to \infty}}
\end{equation*}

\noindent\colorbox{orange}{
  \begin{minipage}{1.0\linewidth}
    \subsection*{Richiesta (c)}
    Per ampiezza di $V_S=1~\text{V}$ e noto che il livello di saturazione negativa del TL081 sia pari
    a $-3.5~\text{V}$ nelle attuali condizioni di alimentazione, calcolare il valore massimo del carico
    tale che il circuito si mantenga lineare e verificare che esso corrisponda approssimativamente a
    quanto osservato al punto 2b.
  \end{minipage}
}
\subsection*{Risposta (c)}
\begin{equation*}
  V_{OUT} = R_Lg_mV_S = -11\frac{R_L}{R}V_S
\end{equation*}
\begin{equation*}
  -3.5~\text{V} = -11\frac{R_L}{10^4}~\text{V} \implies \colorbox{blu}{\boxed{R_L = \frac{35}{11}~\text{k}\Omega}}
\end{equation*}
